\documentclass[12pt,a4paper]{article}
\usepackage[a4paper,margin=18mm]{geometry}
\usepackage[T2A]{fontenc}
\usepackage[utf8]{inputenc}
\usepackage[russian]{babel}
\usepackage{amsmath,amssymb,mathtools}
\usepackage{array,booktabs,longtable}
\usepackage{xcolor}
\usepackage{hyperref}
\usepackage{tikz}

\hypersetup{
  colorlinks=true,
  linkcolor=blue!55!black,
  urlcolor=blue!55!black
}

\setlength{\parindent}{0pt}
\setlength{\parskip}{6pt}
\setlength{\tabcolsep}{3pt}
\renewcommand{\arraystretch}{1.18}

\begin{document}

\pagenumbering{roman}
\begin{titlepage}
\centering
\vspace*{0.27\textheight}

{\LARGE\bfseries Ревью домашней работы\par}

\vspace{1.6cm}

{\large Автор: Лёвин Артём Александрович\par}

\vspace{0.8cm}

{\large 10 сентября 2026 года\par}

\vfill

\begin{tikzpicture}
  \draw[blue!45!black, line width=0.8pt]
    (0,0) .. controls (1.5,0.45) and (3.5,-0.45) .. (5,0)
          .. controls (6.5,0.45) and (8.5,-0.45) .. (10,0);
\end{tikzpicture}

\end{titlepage}

\pagenumbering{arabic}
\newpage
\section*{Сводная таблица}
\small

\begin{tabular}{@{}
>{\raggedright\arraybackslash}p{0.10\linewidth}
>{\raggedright\arraybackslash}p{0.17\linewidth}
>{\raggedright\arraybackslash}p{0.27\linewidth}
>{\raggedright\arraybackslash}p{0.17\linewidth}
>{\raggedright\arraybackslash}p{0.19\linewidth}
@{}}
\toprule
№ задания & Тема & Суть ошибки & Ваш ответ & Правильный ответ \\
\midrule
\hyperref[task:7]{7} &
Упрощение радикальных выражений &
Решение остановлено до финального упрощения; ответ не указан. &
не указан &
$5$ \\
\addlinespace
\hyperref[task:8]{8} &
Сравнение рационального и иррационального чисел &
Верное сравнение получено с помощью десятичных приближений, хотя условие прямо запрещает этот способ. &
$\dfrac{\sqrt2}{3}>\dfrac7{15}$ &
$\dfrac7{15}<\dfrac{\sqrt2}{3}$ \\
\addlinespace
\hyperref[task:10]{10} &
Формулы сокращённого умножения и радикалы &
При раскрытии квадрата суммы вместо $3^2=9$ использовано число $6$. &
$\dfrac{91}{30}$ &
$\dfrac{10}{3}$ \\
\bottomrule
\end{tabular}

\normalsize

\newpage
\null
\vfill
\begin{center}
{\LARGE\bfseries Разбор ошибок}
\end{center}
\vfill
\null

\newpage
\phantomsection
\label{task:7}
\section*{Задание №7}

\textbf{Текст задания.}

Упростите выражение
\[
\frac{\sqrt{45}+\sqrt{20}}{\sqrt5}.
\]

\textbf{Ваше решение.}

В работе выполнены преобразования
\[
\frac{\sqrt{45}+\sqrt{20}}{\sqrt5}
=
\frac{\sqrt5\cdot\sqrt9+\sqrt5\cdot\sqrt4}{\sqrt5}
=
\frac{3\sqrt5+2\sqrt5}{\sqrt5}.
\]
На этом вычисление заканчивается.

\textbf{Ваш ответ.}

Не указан.

\textcolor{red}{Ошибка:} решение не доведено до окончательного упрощения и ответа.

\textbf{Где именно возникает ошибка.}

Последний корректный шаг:
\[
\frac{3\sqrt5+2\sqrt5}{\sqrt5}.
\]
Этот шаг верен: $\sqrt{45}=3\sqrt5$ и $\sqrt{20}=2\sqrt5$.

Первого неверного равенства здесь нет. Проблема возникает после этого шага: преобразование остановлено, хотя выражение ещё можно и нужно упростить. Поэтому решение является неполным.

\textbf{Почему это неверно.}

В задании требуется \emph{упростить выражение}. Запись
\[
\frac{3\sqrt5+2\sqrt5}{\sqrt5}
\]
ещё не является окончательно упрощённым результатом. В числителе находятся подобные слагаемые:
\[
3\sqrt5+2\sqrt5=(3+2)\sqrt5=5\sqrt5.
\]
После этого общий множитель $\sqrt5$ в числителе и знаменателе сокращается. Сокращение допустимо, потому что $\sqrt5>0$, то есть знаменатель не равен нулю.

\textcolor{blue!60!black}{Ключевая идея:} сначала привести корни к виду с одинаковым подкоренным выражением, затем сложить коэффициенты перед $\sqrt5$ и сократить общий множитель.

\textbf{Исправление от последнего правильного шага.}

Продолжим ровно с той строки, на которой решение остановилось:
\[
\frac{3\sqrt5+2\sqrt5}{\sqrt5}
=
\frac{5\sqrt5}{\sqrt5}
=
5.
\]

\textbf{Полное правильное решение.}

Разложим подкоренные числа так, чтобы выделить полный квадрат:
\[
\sqrt{45}=\sqrt{9\cdot5}=3\sqrt5,
\qquad
\sqrt{20}=\sqrt{4\cdot5}=2\sqrt5.
\]
Тогда
\[
\frac{\sqrt{45}+\sqrt{20}}{\sqrt5}
=
\frac{3\sqrt5+2\sqrt5}{\sqrt5}.
\]
Складываем подобные слагаемые в числителе:
\[
3\sqrt5+2\sqrt5=5\sqrt5.
\]
Получаем
\[
\frac{5\sqrt5}{\sqrt5}=5.
\]

\textbf{Проверка.}

После преобразования числитель равен $5\sqrt5$, поэтому
\[
\frac{5\sqrt5}{\sqrt5}=5
\]
при $\sqrt5\neq0$. Следовательно, сокращение выполнено корректно.

\textcolor{teal!60!black}{Итог:} выражение упрощается до
\[
\boxed{5}.
\]

\textbf{Задача на закрепление.}

Упростите выражение
\[
\frac{\sqrt{80}+\sqrt{45}}{\sqrt5}.
\]

\newpage
\phantomsection
\label{task:8}
\section*{Задание №8}

\textbf{Текст задания.}

Сравните без использования десятичных приближений:
\[
\frac7{15}
\qquad\text{и}\qquad
\frac{\sqrt2}{3}.
\]

\textbf{Ваше решение.}

В работе использованы приближённые десятичные значения:
\[
\frac7{15}\approx0{,}46,
\qquad
\sqrt2\approx1{,}41,
\qquad
\frac{\sqrt2}{3}\approx0{,}47,
\]
после чего сделан вывод
\[
\frac{\sqrt2}{3}>\frac7{15}.
\]

\textbf{Ваш ответ.}

\[
\frac{\sqrt2}{3}>\frac7{15}.
\]

\textcolor{red}{Ошибка:} использован способ, прямо запрещённый условием задачи.

\textbf{Где именно возникает ошибка.}

До начала вычислений исходные числа записаны верно. Первый шаг, который не соответствует условию:
\[
\frac7{15}\approx0{,}46.
\]
Далее также используется приближение
\[
\sqrt2\approx1{,}41.
\]
Сам итоговый знак сравнения получен верный, но требование ``без использования десятичных приближений'' не выполнено.

\textbf{Почему это неверно.}

В этой задаче проверяется не только результат сравнения, но и умение сравнивать числа \emph{точно}. Десятичные записи вида $0{,}46$ и $1{,}41$ заменяют исходные числа приближёнными значениями. Это противоречит прямому ограничению в условии.

Кроме того, при сравнении близких чисел грубое округление вообще может изменить или скрыть правильный порядок. Поэтому здесь нужен точный способ, сохраняющий исходные значения.

Оба сравниваемых числа положительны:
\[
\frac7{15}>0,
\qquad
\frac{\sqrt2}{3}>0.
\]
Для положительных чисел можно сравнивать квадраты: если $a>0$ и $b>0$, то
\[
a<b \quad\Longleftrightarrow\quad a^2<b^2.
\]

\textcolor{blue!60!black}{Ключевая идея:} избавиться от корня точным преобразованием --- сравнить квадраты двух положительных чисел.

\textbf{Исправление от последнего правильного шага.}

Вместо перехода к десятичным приближениям сразу сравним квадраты:
\[
\left(\frac7{15}\right)^2=\frac{49}{225},
\]
а
\[
\left(\frac{\sqrt2}{3}\right)^2
=
\frac2{9}
=
\frac{50}{225}.
\]
Так как
\[
\frac{49}{225}<\frac{50}{225},
\]
то, поскольку оба исходных числа положительны,
\[
\frac7{15}<\frac{\sqrt2}{3}.
\]

\textbf{Полное правильное решение.}

Имеем
\[
\frac7{15}>0,
\qquad
\frac{\sqrt2}{3}>0.
\]
Следовательно, можно сравнить квадраты:
\[
\left(\frac7{15}\right)^2
=
\frac{49}{225}.
\]
Для второго числа:
\[
\left(\frac{\sqrt2}{3}\right)^2
=
\frac{(\sqrt2)^2}{3^2}
=
\frac2{9}.
\]
Приведём дробь $\frac2{9}$ к знаменателю $225$:
\[
\frac2{9}=\frac{2\cdot25}{9\cdot25}=\frac{50}{225}.
\]
Теперь сравнение точное:
\[
\frac{49}{225}<\frac{50}{225}.
\]
Значит,
\[
\frac7{15}<\frac{\sqrt2}{3}.
\]

\textbf{Проверка.}

Можно проверить результат другим точным способом. Умножим оба положительных числа на $15$:
\[
\frac7{15}<\frac{\sqrt2}{3}
\quad\Longleftrightarrow\quad
7<5\sqrt2.
\]
Обе части положительны, поэтому возводим их в квадрат:
\[
49<50.
\]
Неравенство верно, значит и исходное сравнение верно.

\textcolor{teal!60!black}{Итог:}
\[
\boxed{\frac7{15}<\frac{\sqrt2}{3}}.
\]

\textbf{Задача на закрепление.}

Сравните без использования десятичных приближений:
\[
\frac5{12}
\qquad\text{и}\qquad
\frac{\sqrt3}{4}.
\]

\newpage
\phantomsection
\label{task:10}
\section*{Задание №10}

\textbf{Текст задания.}

Найдите значение:
\[
\frac{(\sqrt{14}+\sqrt6)(\sqrt{14}-\sqrt6)}
{\sqrt2(\sqrt8+\sqrt2)}
+
\frac{(\sqrt{11}+3)^2-6\sqrt{11}}{10}.
\]

\textbf{Ваше решение.}

В работе первая дробь преобразована так:
\[
\frac{14-6}{\sqrt2\cdot\sqrt8+2}.
\]
Во второй дроби записано:
\[
\frac{11+6\sqrt{11}+6-6\sqrt{11}}{10}
=
\frac{17}{10}.
\]
Далее получено
\[
\frac43+\frac{17}{10}
=
\frac{91}{30}.
\]

\textbf{Ваш ответ.}

\[
\frac{91}{30}.
\]

\textcolor{red}{Ошибка:} при раскрытии квадрата суммы неверно вычислено $3^2$.

\textbf{Где именно возникает ошибка.}

Последний корректный шаг во второй части выражения:
\[
\frac{(\sqrt{11}+3)^2-6\sqrt{11}}{10}.
\]
Первый неверный переход:
\[
(\sqrt{11}+3)^2
\longrightarrow
11+6\sqrt{11}+6.
\]
Средний член $6\sqrt{11}$ найден верно, но последний член должен быть
\[
3^2=9,
\]
а не $6$.

\textbf{Почему этот шаг неверен.}

Используется формула квадрата суммы:
\[
(a+b)^2=a^2+2ab+b^2.
\]
Здесь
\[
a=\sqrt{11},
\qquad
b=3.
\]
Поэтому:
\[
a^2=(\sqrt{11})^2=11,
\]
\[
2ab=2\cdot\sqrt{11}\cdot3=6\sqrt{11},
\]
\[
b^2=3^2=9.
\]
Число $6$ возникает в \emph{среднем} члене $2ab$ как произведение $2\cdot3$. Оно не может одновременно заменять квадрат числа $3$.

Следовательно,
\[
(\sqrt{11}+3)^2
=
11+6\sqrt{11}+9,
\]
а не
\[
11+6\sqrt{11}+6.
\]

Из-за этой ошибки числитель второй дроби получился равным $17$ вместо $20$. Поэтому вторая дробь стала $\frac{17}{10}$ вместо $2$, и итоговое значение всего выражения уменьшилось на
\[
2-\frac{17}{10}=\frac3{10}.
\]

\textcolor{blue!60!black}{Ключевая идея:} при применении формулы $(a+b)^2$ отдельно вычислять три части: $a^2$, $2ab$ и $b^2$, не смешивая коэффициент среднего члена с квадратом второго слагаемого.

\textbf{Исправление от последнего правильного шага.}

Исправим только ошибочную вторую дробь:
\[
(\sqrt{11}+3)^2
=
11+6\sqrt{11}+9.
\]
Тогда
\[
\frac{(\sqrt{11}+3)^2-6\sqrt{11}}{10}
=
\frac{11+6\sqrt{11}+9-6\sqrt{11}}{10}
=
\frac{20}{10}
=
2.
\]
Первая дробь в работе упрощалась в правильном направлении:
\[
\frac{14-6}{\sqrt2\cdot\sqrt8+2}
=
\frac8{\sqrt{16}+2}
=
\frac8{4+2}
=
\frac43.
\]
Следовательно,
\[
\frac43+2
=
\frac43+\frac63
=
\frac{10}{3}.
\]

\textbf{Полное правильное решение.}

Сначала упростим первую дробь. В числителе применяем формулу разности квадратов:
\[
(\sqrt{14}+\sqrt6)(\sqrt{14}-\sqrt6)
=
(\sqrt{14})^2-(\sqrt6)^2
=
14-6
=
8.
\]
В знаменателе раскрываем скобки:
\[
\sqrt2(\sqrt8+\sqrt2)
=
\sqrt2\cdot\sqrt8+\sqrt2\cdot\sqrt2.
\]
Получаем
\[
\sqrt2\cdot\sqrt8=\sqrt{16}=4,
\qquad
\sqrt2\cdot\sqrt2=2.
\]
Значит,
\[
\sqrt2(\sqrt8+\sqrt2)=4+2=6,
\]
и первая дробь равна
\[
\frac86=\frac43.
\]

Теперь упростим вторую дробь:
\[
(\sqrt{11}+3)^2
=
(\sqrt{11})^2+2\cdot\sqrt{11}\cdot3+3^2.
\]
То есть
\[
(\sqrt{11}+3)^2
=
11+6\sqrt{11}+9
=
20+6\sqrt{11}.
\]
Вычитаем $6\sqrt{11}$:
\[
20+6\sqrt{11}-6\sqrt{11}=20.
\]
Следовательно,
\[
\frac{(\sqrt{11}+3)^2-6\sqrt{11}}{10}
=
\frac{20}{10}
=
2.
\]
Складываем результаты:
\[
\frac43+2
=
\frac43+\frac63
=
\frac{10}{3}.
\]

\textbf{Проверка.}

Проверим обе части отдельно:
\[
(\sqrt{14}+\sqrt6)(\sqrt{14}-\sqrt6)=14-6=8,
\]
а знаменатель первой дроби равен
\[
\sqrt2(\sqrt8+\sqrt2)=4+2=6.
\]
Поэтому первая часть действительно равна
\[
\frac86=\frac43.
\]

Для второй части радикальные слагаемые взаимно уничтожаются:
\[
(\sqrt{11}+3)^2-6\sqrt{11}
=
11+6\sqrt{11}+9-6\sqrt{11}
=
20,
\]
значит, вторая часть равна $2$.

Итак,
\[
\frac43+2=\frac{10}{3}.
\]

\textcolor{teal!60!black}{Итог:}
\[
\boxed{\frac{10}{3}}.
\]

\textbf{Задача на закрепление.}

Найдите значение:
\[
\frac{(\sqrt{15}+\sqrt3)(\sqrt{15}-\sqrt3)}
{\sqrt3(\sqrt{12}+\sqrt3)}
+
\frac{(\sqrt7+4)^2-8\sqrt7}{23}.
\]

\vfill
\begin{center}
\small Лёвин Артём Александрович эксклюзивно для Ярослава Верещагина
\end{center}

\end{document}
